% !TeX program = pdflatex
% Two-page English extended abstract. The original Chinese version is kept in
% cohortround_intro.tex.
\documentclass[conference]{IEEEtran}

\usepackage{amsmath,amssymb,bm}
\usepackage{booktabs}
\usepackage{cite}
\usepackage{url}

\setlength{\textfloatsep}{5pt plus 1pt minus 1pt}
\setlength{\floatsep}{4pt plus 1pt minus 1pt}
\setlength{\abovecaptionskip}{2pt}
\setlength{\belowcaptionskip}{-2pt}
\setlength{\abovedisplayskip}{3pt plus 1pt minus 1pt}
\setlength{\belowdisplayskip}{3pt plus 1pt minus 1pt}

\newcommand{\vect}[1]{\bm{#1}}
\newcommand{\diag}{\operatorname{diag}}
\newcommand{\round}{\operatorname{round}}

\title{CohortRound: Sensitivity-Guided Joint Rounding\\
for Fixed-Latency FPGA DSP}

\author{\IEEEauthorblockN{Anonymous Authors}
\IEEEauthorblockA{Three-Page Research Proposal for DAC}}

\begin{document}
\maketitle

\begin{abstract}
Fixed-point compilers normally choose word lengths and round every signal
independently. This misses a hardware-exploitable degree of freedom: errors
introduced on reconvergent DSP paths can cancel at downstream outputs.
We propose \emph{CohortRound}, a compiler and FPGA mechanism that groups a few
physically local, cycle-aligned quantizers and jointly selects their floor/ceil
directions. A bounded-depth selector reads only short discarded-fraction
prefixes and minimizes sensitivity-weighted group discrepancy. The compiler
co-optimizes cohorts, word lengths, alignment, and selector cost, and rejects
cohorts whose control or routing overhead exceeds their datapath savings. An
initial four-term prototype exhaustively validates all 4,096 selector states
and passes bit-true compiler and RTL tests. Post-route FPGA cost remains the
decisive open gate.
\end{abstract}

\begin{IEEEkeywords}
FPGA, DSP, fixed point, word-length optimization, joint rounding, compiler
\end{IEEEkeywords}

\section{Motivation and Research Question}

FPGA word lengths directly affect LUTs, registers, DSP packing, routing, and
energy. Existing word-length optimization (WLO) can assign per-variable formats
under an output-error constraint~\cite{winandy2025}, and static rounding-mode
selection can reduce implementation cost~\cite{menard2010}. Nevertheless, the
runtime decision at each quantizer is still usually local: round-to-nearest
(RTN) minimizes each error separately, not their downstream combination.

Consider quantizers whose outputs later reconverge in a butterfly, FIR/filter
bank, matrix--vector product, or convolution reduction. Independent small
errors may add with the same sign. Conversely, deliberately choosing a locally
non-nearest grid point at one site can cancel errors from other sites. The new
question is therefore: \emph{can a compiler synthesize small, fixed-latency
rounding cohorts whose controlled error correlation permits narrower FPGA
datapaths at the same end-to-end error bound?}

A four-term example illustrates why this differs from local RTN. Suppose the
discarded fractions are $(0.49,0.49,0.49,0.49)$. Independent RTN rounds every
term down and loses $1.96$ output units. A joint policy rounds exactly two terms
up, producing a group error of only $0.04$. Individual errors become larger at
two sites, but the quantity consumed by the DSP reduction becomes much more
accurate. This opportunity is common whenever separately quantized values later
reconverge. It is absent when errors reach unrelated outputs, have incompatible
timing, or require excessive wiring; identifying that boundary is part of the
compiler problem rather than an assumption of the method.

Prior filter work shapes roundoff noise by selecting static widths
\cite{constantinides2000}, while correlation models estimate noise already
created by a datapath~\cite{naud2012}. Controlled/dependent rounding provides
joint-rounding foundations~\cite{coxernst1982,gandhi2006}, and DiscQuant applies
discrepancy ideas to offline neural-network weights~\cite{chee2025}. Stochastic
rounding hardware instead makes local PRNG-driven decisions~\cite{mikaitis2021}.
None directly synthesizes a sensitivity-aware, per-sample joint selector under
FPGA pipeline, logic-depth, and physical-locality constraints.

\section{CohortRound Mechanism}

\subsection{Joint discrepancy}
For site $i$ with step $\Delta_i$, write
$x_i=\Delta_i(n_i+a_i)$, where $n_i=\lfloor x_i/\Delta_i\rfloor$ and
$a_i\in[0,1)$. A legal floor/ceil decision $r_i\in\{0,1\}$ introduces
\begin{equation}
 \delta_i=\Delta_i(r_i-a_i).
 \label{eq:local}
\end{equation}
For a cohort $C$ in an affine downstream region, its monitored-output error is
\begin{equation}
 \vect e_C=A_CD_C(\vect r_C-\vect a_C),\quad
 D_C=\diag(\Delta_i),
 \label{eq:vector}
\end{equation}
where columns of $A_C$ are downstream sensitivities. Instead of independent
RTN, an ideal cohort selects
\begin{equation}
 \vect r_C^*=\arg\min_{\vect r\in\{0,1\}^{|C|}}
 \|WA_CD_C(\vect r-\vect a_C)\|_\infty .
 \label{eq:oracle}
\end{equation}
The compiler initially admits exact affine regions; nonlinear regions require a
sound Jacobian/remainder envelope rather than sampled sensitivities.

\subsection{Prefix selector and certificate}
Sending every discarded bit to a shared unit would erase the saving. The
selector therefore observes only $b$ leading residue bits
$p_i=\lfloor2^ba_i\rfloor$ and implements a small truth table
$T_C(\vect p)\in\{0,1\}^{|C|}$. Its deterministic certificate is
\begin{equation}
 \Gamma_C=\max_{\vect p}\sup_{\vect a\in\mathcal B(\vect p)}
 \|WA_CD_C(T_C(\vect p)-\vect a)\|_\infty .
 \label{eq:cert}
\end{equation}
For $k$ equal-sensitivity terms with common step $\Delta$, choose
$m=\round(\sum_i a_i)$ and any binary vector with $\sum_i r_i=m$. Then the
group error is at most $\Delta/2$, whereas independent RTN can approach
$k\Delta/2$. With prefix midpoints
$\tilde a_i=(p_i+1/2)2^{-b}$, the bound becomes
\begin{equation}
 |e_C|\le \Delta\left(\frac12+\frac{k}{2^{b+1}}\right).
 \label{eq:prefixbound}
\end{equation}
Thus the implemented $k=4,b=3$ selector has a $0.75\Delta$ bound while using
only twelve residue bits and a three-bit correction count.

For unequal sensitivities, the identity of the rounded-up lanes matters, not
only their count. HCDR enumerates each prefix cell and solves the robust local
problem
\begin{equation}
 T_C(\vect p)=\arg\min_{\vect r\in\{0,1\}^{|C|}}
 \sup_{\vect a\in\mathcal B(\vect p)}
 \|WA_CD_C(\vect r-\vect a)\|_\infty .
 \label{eq:robustcell}
\end{equation}
The resulting mapping is minimized into LUTs, comparators, and a small adder
tree. Prefix width is itself a design variable: increasing $b$ tightens the
certificate but adds source fanout, wires, and selector inputs. Candidate
selectors that miss the cycle budget are either pipelined with explicitly
charged alignment registers or rejected.

\subsection{Hardware-Constrained Discrepancy Rounding}
Our compiler algorithm, HCDR, performs four steps. First, it extracts site
scales, sensitivities, lifetimes, pipeline stages, and placement regions from a
scheduled DSP graph. Second, it constructs a candidate-cohort hypergraph and
uses a software joint-rounding oracle to estimate attainable bit savings.
Third, exact truth-table/MILP search synthesizes small prefix selectors subject
to depth, fanout, and span limits; larger groups use a sensitivity-aware
heuristic measured against the exact small-instance oracle. Finally, a
hardware-aware set-packing pass selects non-overlapping cohorts jointly with
word lengths and schedule alignment:
\begin{align}
 \min_{\vect f,\vect z,\{T_C\},S}\quad &
 C_{\rm data}(\vect f,S)+\sum_Cz_C(C_{\rm sel}+C_{\rm route}) \notag\\
 \text{s.t.}\quad &\Gamma_{\rm single}+\sum_Cz_C\Gamma_C\le\epsilon,
 \quad \mathrm{II}\le\mathrm{II}_{\max}. \label{eq:global}
\end{align}
Timing, alignment, overflow, and one-cohort-per-site constraints are also
enforced. This explicit cost model lets the compiler reject unprofitable
cohorts rather than assuming joint rounding is always beneficial.

\subsection{What is genuinely new}
Joint rounding by itself is not new. The proposed contribution is the combined
compiler/hardware object
\begin{equation}
 \textit{cohort}=(C,A_C,D_C,b,T_C,\Gamma_C,\textit{placement window}),
 \label{eq:object}
\end{equation}
which connects compile-time sensitivities and physical constraints to a
per-sample hardware policy. WLO changes static formats but not the correlated
runtime decisions. Dependent rounding supplies mathematical tools but does not
map dynamic DSP intermediates onto a bounded-latency circuit. Correlation
analysis predicts existing noise but does not synthesize its signs. DiscQuant
uses discrepancy to round static model weights offline. Stochastic rounding is
runtime, but each decision is PRNG-driven rather than jointly conditioned on
downstream sensitivities.

The distinction from delayed rounding is also essential. If all cohort members
feed one immediate sum with equal sensitivity, retaining guard bits and rounding
once is a strong baseline and may be preferable. CohortRound becomes structurally
distinct when values are quantized before traversing different branches,
buffers, or operators and later affect one or more outputs with unequal signed
sensitivities. In such graphs, moving one common rounding point to the end is not
equivalent without preserving wide data across the intervening region. The
evaluation therefore includes both direct reductions and branched/reconvergent
graphs; success on the former alone is insufficient.

\section{Prototype Evidence and DAC Evaluation}

We added a \texttt{cohort\_mac\_reduce} intrinsic, a fail-closed canonical
dot-product recognizer, bit-true interpreter semantics, Chisel lowering, and
three matched RTL designs. All use 24 signed $\mathrm{Fixed}[16,14]$ terms, the
same digit-serial multiplier and control schedule, and a
$\mathrm{Fixed}[20,13]$ output. Only the rounding/reduction policy changes:
CohortRound, independent per-product RNE, or full-precision accumulation
followed by one RNE.

The current target forms exact products sequentially. For every group of four,
it arithmetic-shifts each signed product to a 17-bit floor, collects the three
most significant discarded bits, and accumulates the floors in 19 bits. The
selector converts the four prefixes into a correction in $[0,4]$, which is
added to a 23-bit global accumulator. Negative products use the same
floor--residue decomposition, so no sign-specific approximation is hidden in
the software model. Two matched controls preserve input/output formats,
operation count, multiplier, handshake, and schedule: one independently rounds
every product, while the other accumulates in 38 bits and rounds once.

\begin{table}[t]
\centering
\caption{Completed pre-Vivado gates.}
\label{tab:prototype}
\footnotesize
\begin{tabular}{@{}lrl@{}}
\toprule
Gate & Result & Coverage \\
\midrule
Prefix oracle & 4,096/4,096 & all $8^4$ states \\
Bit-true datapath & 2,000/2,000 & random signed MACs \\
Compiler tests & 7/7 & checker, lowering, fail-closed \\
Chisel/Verilator & 5/5 & target and two baselines \\
RTL lint & 3/3 & all generated tops \\
\bottomrule
\end{tabular}
\end{table}

The lowering reduces the global accumulator from 38 to 23 bits and uses a
19-bit four-term group accumulator. However, the present serial multiplier
still forms a full 32-bit product; therefore these results establish numerical
and RTL feasibility, not a large total-area saving.

The tests deliberately cover more than average error. The Python selector
exhausts all prefix states, including ties; an independently written integer
oracle checks negative operands, extreme products, output wraparound, and 2,000
random vectors. RTL tests repeat exhaustive selector checking, verify the
one-cycle valid/bubble protocol, and cross-check the complete 24-term machine
and both baselines. Unsupported cohort sizes and noncanonical kernels fail
closed instead of silently falling back to different arithmetic.

The full evaluation will cover FIR/polyphase filters, FFT butterflies,
matrix--vector/beamforming, and convolution reductions, with negative controls
that have no local alignable cohort. Baselines include truncation, independent
RTN and stochastic rounding, static mode optimization~\cite{menard2010}, and
width-only WLO~\cite{winandy2025}. We will report post-route LUT/FF/DSP/BRAM,
Fmax, II, latency, energy/sample, routing/alignment overhead, error certificates,
and exact-oracle gap. The idea passes only if at least three basic DSP kernels
save at least one fractional bit at the same deterministic error bound, retain
II, lose less than 5\% Fmax, and achieve net post-route area or energy benefit.
If the shared multiplier dominates, a truncated-product/high-product-plus-prefix
lowering is required; changing the device or timing target is not evidence.

To isolate causality, the study uses four ablations: (i) identical widths with
independent decisions, measuring the value of correlation; (ii) oracle cohorts
without locality constraints, measuring the physical feasibility gap; (iii)
identical cohorts with full residues versus short prefixes, measuring the
information--routing trade-off; and (iv) compiler-selected versus random or
cardinality-matched cohorts, measuring the value of sensitivity analysis. Each
design is synthesized under the same part, clock, II, I/O, and multiplier
architecture. Power claims require workload-derived SAIF/VCD activity rather
than vectorless estimates.

Three failure modes are reported explicitly. First, many kernels may expose
only isolated quantizers or sensitivities that do not permit cancellation.
Second, the selector and alignment network may consume more resources than one
saved fractional bit. Third, a serial architecture may be multiplier-dominated,
making reduction-width savings invisible at top level. The last case motivates
a high-product-plus-prefix multiplier, but that extension must be evaluated as
a separate mechanism and ablation rather than credited to cohort rounding.

\subsection{Compiler integration and supported scope}
The front end first lowers arithmetic casts and explicit quantization sites into
an IR that records source/destination formats, discarded-bit ranges, and
overflow semantics. A backward sensitivity pass annotates each site with its
transfer vector to selected outputs. Scheduling then supplies availability
cycles and lifetimes, while a placement proxy supplies candidate regions and
wire-distance estimates. HCDR consumes these annotations and emits ordinary
fixed-width operators plus a selector module, residue-prefix connections, and
any required delay registers. The bit-true interpreter evaluates exactly the
same prefix policy used by RTL, preventing a floating-point software oracle from
silently defining different tie or signed-shift behavior.

The first compiler version deliberately targets real-valued affine regions and
acyclic reductions. Saturation and overflow are orthogonal: range analysis must
prove no overflow or model the selected wrap/saturate semantics before a
rounding certificate is valid. Feedback filters require an invariant over both
data state and rounding state and are therefore a separate extension. Likewise,
the stateless selector provides a per-invocation deterministic discrepancy
bound; it does not claim probabilistic unbiasedness. A future debt controller
may reduce long-term drift only if finite-state reachability proves its debt is
bounded for every legal input sequence.

The expected deliverables are consequently concrete: a cohort-aware numerical
IR and discovery pass; an exact small-cohort robust selector algorithm with a
scalable heuristic; machine-checkable error certificates that include unseen
residue suffixes; and an RTL backend with post-route cost feedback. Releasing
rejected candidates and their causes is also important: it reveals whether the
limiting factor is numerical structure, schedule alignment, selector logic, or
routing, and prevents reporting only favorable kernels.

\section{Conclusion}
CohortRound turns runtime correlation among quantization errors into a
compiler-visible FPGA design variable. Its novelty is not discrepancy rounding
alone, but sensitivity-aware cohort discovery, fixed-latency prefix-selector
synthesis, deterministic certification, and physical-cost-aware extraction.
The completed prototype supports the numerical and implementation mechanism;
matched Vivado post-route results remain the decisive test of DAC-level impact.

\bibliographystyle{IEEEtran}
\bibliography{references}

\end{document}
